% MSc dissertation example file, February 2022
%
% Leave one of the documentclass lines uncommented to match your degree.
% You may remove the logo option if it causes problems.
% Do not change any other options.
% \documentclass[logo,msc,adi]{infthesis}     % Adv Design Inf
% \documentclass[logo,msc,ai]{infthesis}      % AI
% \documentclass[logo,msc,cogsci]{infthesis}  % Cognitive Sci
% \documentclass[logo,msc,cs]{infthesis}      % Computer Sci
% \documentclass[logo,msc,cyber]{infthesis}   % Cyber Sec
% \documentclass[logo,msc,datasci]{infthesis} % Data Sci
% \documentclass[logo,msc,di]{infthesis}      % Design Inf
% \documentclass[logo,msc,dsti]{infthesis}    % Data Sci TI
% \documentclass[logo,msc,inf]{infthesis}     % Informatics
\documentclass[logo,msc,ai]{infthesis}           % degree unspecified, do not change except to add your degree
%%%%%%%%%%%%%%%%%%%%%%%%
% Understand any problems and seek approval before assuming it's ok to remove ugcheck.
% \usepackage{msccheck}
\usepackage{Original/msccheck} %LAST SUBMISION CHANGE TO ABOVE AND MOVE FILE BACK OUT OF ORIGINAL

% Include any packages you need below, but don't include any that change the page
% layout or style of the dissertation. By including the ugcheck package above,
% you should catch most accidental changes of page layout though.

\usepackage{microtype} % recommended, but you can remove if it causes problems
\usepackage{xcolor} %remove later this is used for helping edit and draft
\usepackage{subcaption} %subcaptions for plots and figures next to each other. 
\usepackage{amsmath}  %Check we are allowed
\usepackage{amssymb} %Check we are allowed
\usepackage{amsfonts} %Check we are allowed
\usepackage{booktabs} %Check we are allowed
\usepackage{graphicx} %Check we are allowed
\usepackage{tabularx}
\usepackage{float}


%%% Colour commands - delete for final draft 
% - Editing Macros (Hierarchy of Importance) -
% 1. FACT CHECK (Critical): Red commands immediate attention for potential inaccuracies.
\newcommand{\fc}[1]{\textcolor{red}{[FactC: #1]}} 
% 2. CITE (High): Orange highlights missing evidence or attribution.
\newcommand{\ct}[1]{\textcolor{orange}{[CITE: #1]}} 
% 3. APPENDIX (Medium): Blue indicates structural moves (shifting text out of the main body).
\newcommand{\app}[1]{\textcolor{blue}{[APP: #1]}} 
% 4. SMALL REWORD (Low): Teal is for minor stylistic edits, mirroring your friend's setup.
\newcommand{\sr}[1]{\textcolor{teal}{[Sreword: #1]}}
% 5. FIGURE (Low): Purple is for figure-related notes (refer to figure, create figure, etc.).
\newcommand{\fig}[1]{\textcolor{purple}{[FIG: #1]}}
% 6. AI redraft (Low): Ai marker suggestions and corrections
\newcommand{\ai}[1]{\textcolor{brown}{[AI: #1]}}
% 7. LIMITATION (Low): Move to limiations 
\newcommand{\limit}[1]{\textcolor{cyan}{[LIM: #1]}}

%=======================================================================================
\begin{document}
\begin{preliminary}

\title{Physics-informed prediction and spatial attribution of forest top-height growth using repeated LiDAR surveys}
%TODO: DRAFT DIFFERENT VERSIONS OF THE TITLE AND PICK AGAIN AT THE END

\author{Shreeya Gorasia}

\date{\today}
%=======================================================================================

\abstract{TODO: FILL IN}

%TODO: writing: Comeback to this section when the entire report is done 
%=======================================================================================

\maketitle

\newenvironment{ethics}
   {\begin{frontenv}{Research Ethics Approval}{\LARGE}}
   {\end{frontenv}\newpage}

\begin{ethics}
This project was planned in accordance with the Informatics Research
Ethics policy. It did not involve any aspects that required approval
from the Informatics Research Ethics committee.

\standarddeclaration
\end{ethics}

%TODO: writing: Comeback to this section when the entire report is done 
%=======================================================================================

\begin{acknowledgements}
Any acknowledgements go here.
\end{acknowledgements}

%TODO: writing: Comeback to this section when the entire report is done 
%=======================================================================================

\tableofcontents
\end{preliminary}

%TODO: check: Comeback to this when entire report is done and check it and make it fits in one page
%=======================================================================================
\chapter{Introduction} \label{chap:introduction}

\section{Motivation and problem context}
% Forestry problem and wider importance
A Forest Digital Twin is a wider long-term goal for the forestry sector: a dynamic digital representation of forest systems that can support monitoring, forecasting and management planning ~ a shift from static, periodic assessment toward adaptive, data-driven planning. The aim is to combine repeated observations, simulation, and machine-learning(ML) and artificial intelligence (AI) to represent forests as changing systems \cite{baskent2026Review_C1_SS, buonocore2022Proposal_C1_SS}. This is increasingly vital as managers face climate-driven risks—such as extreme storms, destructive winds, and severe summer heatwaves that trigger wildfires and alter growth patterns. At the same time, global forestry is shifting toward continuous cover management, requiring more flexible stand-level models \cite{ForestResearch2009Evidence_C1_SS}. 

Sitka spruce (\textit{Picea sitchensis}) is Britain's primary commercial species, covering approximately 50\% of managed forests \cite{ForestResearch2026Sitka_C1_SS} and contributing roughly \pounds1.1 billion annually to Scotland's economy \cite{ScottishForestry2026Future_C1_SS, mason2011Sitka_C1_SS}. The study focuses on Aberfoyle in Trossachs, Scotland, a Sitka spruce-dominated region with 21 years of repeated LiDAR coverage. Featuring varied topography with exposed, windy ridges to sheltered valleys, providings variation that causes different growth trajectories in different areas of the forest. 

% Height, LiDAR and shared growth relationship
Top height is a key measure of stand development, commonly used in forestry models to estimate downstream metrics like timber volume and site productivity \cite{ForestResearch2026Sitka_C1_SS}, and now can be accurately captured at scale using airborne LiDAR scans \cite{suarezminguez2008Practical_C1_SS}. For most tree species, including Sitka spruce, height development over time is modelled by the Chapman-Richards equation \cite{pommerening2015Methods_C2_CR}, a sigmoid growth curve, with its parameters fixed to represent a stand's maximum height potential and its speed of growth. A fixed parameter curve only captures the average population trend, but local trajectories can deviate from this baseline, due to factors like terrain, climate, soil, wind exposure, and active management all critically constrain or enhance a stand's actual development \cite{socha2023Higher_C1_SS, cameron2015Building_C1_SS}. This creates two related questions: whether incorporating these local environmental conditions can improve spatial height prediction, and whether they are associated with plots departing from the shared growth relationship.

Flexible ML architectures can capture complex, non-linear interactions across age, stand structure, and local environment. While models like deep neural networks (DNNs) and gradient-boosted decision trees (XGBoost) learn these relationships directly from observation, physics-informed neural networks (PINNs) provide a hybrid approach by embedding domain knowledge directly into the training objective. DNNs are known to struggle with similar growth processes, as its susceptible to overfitting or data sparsity, creating unplausible predictions for example decreasing height over time. The PINN instead uses the CR equations and derivate terms in the loss function to act as regularisation. These soft loss constraints guide predictions toward smooth, plausible growth shapes, while incorporating environmental features allows the model to capture how local conditions alter the growth trajectory. \cite{karpatne2017Theory_C4_NN, pichler2025Inferring_C4_NN}

\paragraph*{Research gap and Contribution}
While repeated airborne laser scanning has been applied to model forest height growth and productivity \cite{janiec2023Development_C1_SS}, limited evidence exists on whether Chapman Richards hybrid model improves predictions of Sitka spruce top height, and using these same models as a tool to more understand the correlations between abiotic factors and forest growth. Many current remote-sensing models do not account for spatial dependence between nearby-identical trees, leading to data leakage and inflated performance estimates \cite{Karasiak2022Spatial_C5_EVAL, ploton2020Spatial_C5_EVAL, lynch2025Digital_C4_NN}. Relying on black-box neural networks without explainable AI limits their use for stakeholders who require interpretable insights to support decision-making.

Therefore the focus of this dissertation targets a small component of the Digital Twin pipeline: testing whether incorporating multi-source environmental data and ecological theory improves spatial height predictions, and using attribution techniques to explain local growth deviations.

\vspace{5 em}
\vspace{-0.5em}

\section{Objectives and Research Questions}

\vspace{-0.5em}
\begin{table}[!htbp]
\small
\setlength{\tabcolsep}{4pt}
\renewcommand{\arraystretch}{0.95}

\noindent\makebox[\linewidth][c]{%
\begin{tabularx}{1.08\linewidth}{
    @{}
    >{\raggedright\arraybackslash}p{0.18\linewidth}
    >{\raggedright\arraybackslash}X
    @{}
}
\noalign{\hrule height 0.9pt}

\textbf{Linked RQ}
&
\textbf{Objective, Sub-objectives \& Own Contribution}
\\

\noalign{\hrule height 0.9pt}

\textbf{All RQs}
\newline
Dataset
\newline
construction
&
Clean and combine repeated forest inventory and LiDAR scan records,
with multi-source geospatial rasters (terrain, wind, climate, and soil) into
a unified modelling dataset.
\\

\hline

\textbf{RQ1}:
\newline
Top-Height
\newline
Prediction
&
\textbf{(1.1)} Compare ML baselines (LR, RF, XGBoost) against deep neural
networks (DNNs) and physics-informed neural networks (PINNs).
\newline
\textbf{(1.2)} Evaluate generalisation under spatial (unseen sites) and
temporal (unseen survey years) holdouts.
\newline
\textbf{(1.3)} Quantify accuracy gains from adding multi-source environmental
predictors.
\newline
\textbf{Contribution:} Comparison ML, DNN and PINN models for top height
prediction under spatial cross-validation.
\\

\hline

\textbf{RQ2}:
\newline
CR-curve
\newline
Deviation
\newline
Analysis
&
\textbf{2.1)} Test whether environmental information reduces the deviations
from the average growth trajectory.
\newline
\textbf{(2.2)} Identify environmental and management conditions associated
with persistent plot-level departures.
\newline
\textbf{Contribution:} Show which specific conditions associated with
persistent departures from the shared growth curve
\\

\hline

\textbf{RQ3}:
\newline
Plot-specific
\newline
Variation
&
\textbf{3.1)} Estimate plot-specific asymptotic height and its departure from
expected growth.
\newline
\textbf{3.2)} Assess the reliability of the estimated plot-level differences.
\newline
\textbf{3.3)} Identify associated environmental and management conditions
using global models.
\newline
\textbf{3.4)} Test whether allowing these associations to vary spatially
improves held-out performance.
\newline
\textbf{Contribution:} Distinguishe real environmental trends from noise
and fitting errors
\\

\noalign{\hrule height 0.9pt}
\end{tabularx}%
}

\label{tab:aims_objectives}
\end{table}
\vspace{-0.8em}

%=======================================================================================
%=======================================================================================
\chapter{Related Work and Background} \label{ch:background}
This chapter reviews tree growth modelling, machine learning and physics-informed prediction, and spatial attribution methods. It concludes by highlighting the research gaps this study addresses.

\section{Top-height growth and environmental variation}

The \textbf{Chapman-Richards (CR)} equatioon is widely used in forestry to model height using a flexible sigmoidal curve \cite{richards1959Flexible_C2_CR,pommerening2015Methods_C2_CR,manso2021Dynamic_C2_CR}:
\begin{equation}
H(a) = y_{\max}\left(1 - e^{-ka}\right)^p,
\label{eq:background-cr}
\end{equation}
where $H(a)$ is dominant top height at stand age $a$, $y_{\max}$ is the fitted asymptotic height ceiling (typically 40--50\,m in plantation settings, distinct from the biological maximum usually 80\,m; \cite{puliti2023Tree_C2_CR}), $k$ controls the approach rate to the asymptote, and $p$ governs the inflection point and trajectory shape.

When observational data span a narrow age range, isolating $y_{\max}$ directly is difficult because different combinations of $(y_{\max}, k, p)$ produce overlapping curves over intermediate ages. Observational records also carry measurement noise, repeated-survey correlation, and stand disturbance that functional form alone cannot resolve \cite{rennolls1995Forest_C2_CR}.  British forestry use this relationship to measure site productivity defined as General Yield Class (GYC): the peak average yearly timber volume growth over a stand's life ($\mathrm{m^3\,ha^{-1}\,yr^{-1}}$) \cite{edwards1981Yield_C1_SS,ForestResearch2026Sitka_C1_SS}. For Sitka spruce, empirical yield tables solve directly for GYC by inverting the Chapman--Richards equation:
\begin{equation}
\mathrm{GYC}_{\mathrm{SS}} \approx \alpha \cdot \frac{H(a)}{\left(1 - e^{-k a}\right)^p} - \beta,
\label{eq:gyc-ss}
\end{equation}
where $H(a)$ and $a$ are observed top height and age, while $\alpha$, $\beta$, $k$, and $p$ are calibrated constants for the species (full parameter values given in Appendix~\ref{app:gyc-parameters}). 

%TODO: METNION IN METHOD NOT BACKGROUND: Because GYC is calculated directly from top height via \eqref{eq:gyc-ss}, using inventory-derived productivity ratings as predictors of top height would introduce target leakage.

\begin{figure}[!htbp]
  \centering
  \includegraphics[width=0.65\textwidth]{figures/fig_background/CR-curves-yeildclass.png}
  \caption{\small{CR growth curves Equation~\eqref{eq:background-cr}, across General Yield Classes (GYC 2--24; $\mathrm{m^3\,ha^{-1}\,yr^{-1}}$) for Sitka spruce in Aberfoyle. Higher yield classes scale toward greater asymptotic heights ($y_{\max}$), inverted by Equation~\eqref{eq:gyc-ss} to infer productivity from height and age. Curves reflect regional averages, with individual stands deviating under local site conditions.}}
  \label{fig:chapman-richards-illustrative}
\end{figure}

\paragraph*{Environmental drivers of height growth.} Site-conditioned and mixed-effects have been used to adjust growth-curve parameters using local \emph{physiographic} (terrain and topography) and \emph{bioclimatic} covariates \cite{scolforo2013Dominant_C2_CR,socha2021Regional_C2_CR}. For Sitka spruce across the UK and Ireland, height growth is primarily shaped by elevation, wind exposure, soil moisture, and nutrient availability, for example in northern UK, GYC drops by $3.2\text{--}4.0\,\mathrm{m^3\,ha^{-1}\,yr^{-1}}$ per 100\,m rise in elevation, driven mainly by cooler temperatures and stronger winds \cite{worrell1990Productivity1_C1_SS}. In Ireland studies show soil nutrient levels explain over half the variation in site productivity, with growth slowing at both soil moisture extremes and wind exposure limiting growth primarily on high-elevation sites \cite{farrelly2011Sitka_C1_SS}.

Mapping forest growth across landscapes often pairs inventory records with geospatial data. In Wales, early UK applications linked Forestry Commission sub-compartment records with terrain and soil layers via PCA and regression to model Sitka spruce yield class \cite{bateman1998Using_C1_SS}. Modern approaches combine airborne LiDAR canopy heights with multi-source environmental rasters, demonstrating that macro-climate sets broad height potential while local terrain and soil drive stand-level variation \cite{fricker2019More_C3_LID}. These same data families supplement the Forest Research plot records in this study. However, model reliability hinges on layer resolution; coarse data-sources like SoilGrids show poor ground agreement ($R^2 = 0.19$; requiring careful local validation \cite{miller2024Assessing_C2_CR}).

\section{Data-driven and process-informed prediction} \label{sec:background-prediction}
ML models learn patterns directly from data without assuming a fixed curve shape \cite{lee2018Machine_C3_LID}. Random Forest averages independently fitted trees through bagging and random feature selection, to reduce variance \cite{breiman2001RandomForest_C4_NN}. Gradient Boosted Decision Trees (XGBoost), build trees sequentially to correct the residual errors of the previous ones \cite{chen2016XGBoost_C4_NN}. Both approaches handle non-linear relationships and feature interactions well without requiring manual parameter curves. However, forest remote-sensing studies often test these models on single-date snapshots; predicting repeated height growth requires models to generalize to unseen forest blocks, across stand ages, and over multi-year survey intervals.

\paragraph*{Neural models of forest growth.}
Deep neural networks (DNNs) learn non-linear interactions through layered hierarchical representations, optimising their weights through backpropagation to minimise a specified training loss. But these do not inherently preserve properties such as monotonic growth or bounded asymptotes. Previous comparisons show that unconstrained neural networks can reproduce broad sigmoid growth patterns while failing to preserve these properties across age ranges \cite{ercanli2023Comparison_C4_NN, karatepe2022Total_C4_NN}.

\textbf{PINNs} and hybrid models incorporate scientific relationships into the training objective \cite{karpatne2017Theory_C4_NN,raissi2019Physics_C4_NN, karniadakis2021Physics_C4_NN,willard2022Integrating_C4_NN}. A common loss formulation is:
\begin{equation}
\mathcal{L}_{\text{total}} = \mathcal{L}_{\text{data}} + \lambda\mathcal{L}_{\text{process}},
\label{eq:pinn-loss-general}
\end{equation}

where $\mathcal{L}_{\text{data}}$ measures prediction error, $\mathcal{L}_{\text{process}}$ measures disagreement with the selected process relationship, and $\lambda$ balances the two objectives.

Process knowledge can enter neural models in several ways
\cite{wesselkamp2024Process_C4_NN}:

\begin{table}[ht]
\centering
\footnotesize
\begin{tabular}{ll}
\hline
\textbf{Integration Strategy} & \textbf{Structural Mechanism} \\
\hline
\emph{Bias correction} & Network learns systematic residuals of process model outputs. \\
\emph{Parallel physics} & Process and neural predictions are combined additively. \\
\emph{Physics regularisation} & Unconstrained neural predictor is penalised for equation disagreement. \\
\emph{Domain adaptation} & Pre-training on process simulations before empirical fine-tuning. \\
\emph{Physics embedding} & Differentiable process models embedded with parameter sub-networks. \\
\hline
\end{tabular}
\label{tab:pinn-strategies}
\end{table}

% Conditional and variable-coefficient PINN frameworks combine physics regularisation with physics embedding, letting environmental covariates drive growth-curve parameters rather than treating them as fixed constants \cite{kovacs2022Conditional_C4_NN,miao2023VCPINN_C4_NN,scolforo2013Dominant_C2_CR}. Frameworks such as Forest-Informed Neural Networks integrate neural networks into more complex differential equation models covering growth, mortality, regeneration, and competition; a single Chapman--Richards curve is a simpler target, for which the ecological equation can instead be anchored with a soft penalty term \cite{pichler2025Inferring_C4_NN}. -- what i would of done if i did both the physics reglarisation over the hard embedding and the keeping terrain out of the raw input rather than cocntancating: what i didnt do is both together. 

Conditional and variable-coefficient PINN frameworks combine physics regularisation with physics embedding, letting environmental covariates drive growth-curve parameters rather than treating them as fixed constants. Frameworks such as Forest-Informed Neural Networks integrate neural networks into more complex differential equation models covering growth, mortality, regeneration, and competition; a single Chapman--Richards curve is a simpler target, for which the ecological equation can instead be anchored with a soft penalty term \cite{pichler2025Inferring_C4_NN}.

Adding process penalties creates a trade-off between fitting the observations and adhering to the Chapman--Richards guidance. Composite loss functions can produce imbalanced gradients, causing one loss component to dominate optimisation, while ill-conditioned loss landscapes can slow or prevent convergence \cite{wang2021Understanding_C4_NN,rathore2024Challenges_C4_NN}. There is also a risk that a strong penalty based on a misspecified population-level Chapman--Richards curve pulls predictions towards an unsuitable average rather than allowing the network to represent local growth differences.


%------------------------------------------------
\section{Spatial Evaluation and Growth Attribution}

A common problem when workign with spatial data is that nearby plots share similar soils, climate, and management histories, and thus models can cheat when using random train--test splits as models memorise geographic locations rather than generalisable biological relationships \cite{ploton2020Spatial_C5_EVAL, karasiak2022Spatial_C5_EVAL}. To properly test how well models transfer studies stress the need to test them on separate locations, times, or both To properly test how well models transfer previous work state the importance of using spatially and temporally disjoint holdouts \cite{meyer2018Improving_C5_EVAL}.

Beyond raw prediction, research has evaluated methods for attributing growth differences to specific environmental drivers. Global regularised models such as Elastic Net balance sparsity and collinearity by combining $L_1$ and $L_2$ penalties, yielding interpretable coefficients under assumptions of linearity and spatial stationarity \cite{zou2005Regularization_C4_NN}. Flexible tree ensembles like XGBoost capture complex non-linearities and interactions, with TreeSHAP decomposing predictions into additive feature contributions \cite{chen2016XGBoost_C4_NN,lundberg2017Unified_C7_XAI}. However, SHAP shows how the model makes decisions rather than true biological cause and effect, and correlated variables can distort how important each feature looks.

Past studies use spatial regression to test whether environmental associations vary across space. Standard Geographically Weighted Regression (GWR) fits a local linear regression at each location, weighting nearby observations more heavily than distant ones \cite{brunsdon1996Geographically_C6_GWR}. Geographically Neural Network-Weighted Regression (GNNWR) replaces the fixed distance kernel with a neural network that learns spatial weights directly from coordinate distances, allowing local coefficients to vary more flexibly \cite{du2020Geographically_C6_GWR}. In China, a GNNWR was used to model fallen trees from terrain and stand variables, outperforming global models and classical GWR while reducing residual spatial autocorrelation \cite{sun2022Local_C6_GWR}. However, 355 of the 443 plots were randomly selected for training and the remaining 88 for validation within the same reserve. The study therefore showed that learned spatial weights captured variation within the sampled reserve, but they did not test transfer to spatially disjoint areas within that reserve or to a separate forest.

%------------------------------------------------
\section{Research gap and study position}
\label{sec:background-gap}

Established top-height models provide simple representations of stand development, yet uniform height-age curves often mask significant local growth variation. Multi-temporal LiDAR data provides extensive spatial coverage of dominant canopy height, but these metrics reflect a combination of biological growth, sensor acquisition differences, and unmeasured plot-level management.

ML has already been used to relate LiDAR-derived canopy height to environmental conditions in other forests \cite{fricker2019More_C3_LID}. However, no study identified in this review combines repeated LiDAR observations with stand records and GIS-derived terrain, wind, climate and soil variables to model Sitka spruce height trajectories at Aberfoyle, where direct field measurements are limited. ML and PINNs provide flexible alternatives to static parametric curves, but adding domain guidance does not inherently guarantee better predictions than an unconstrained DNN or a tuned tree ensemble. Separately, the attribution of plot-level growth-curve departures has not been compared across linear, tree-based and spatially varying models under spatially independent validation, leaving it uncertain whether different model families identify consistent environmental influences.

This dissertation addresses these gaps through the aim and objectives set out in the introduction.


% Explain what previous studies did.
% State the important findings.
% Compare the approaches.
% Identify limitations that lead directly to the dissertation’s research gap.


%=======================================================================================
%=======================================================================================
\chapter{Data} \label{ch:data}
This chapter describes the data used for model training and evaluation, including cohort construction, feature engineering and the main data limitations. Table~\ref{tab:data_summary} summarises the predictor and target variables, their sources and temporal coverage. The rationale for feature selection, data splitting and model evaluation is given in Chapter~\ref{ch:methodology}.

\section{Study area and cohort coverage}
The dataset covers approximately $16\text{~km} \times 12\text{~km}$ of Aberfoyle and contains 795,645 raw plot observations across six survey years (2002--2023). The records contain airborne laser scanning (ALS)-derived variables, including canopy cover and height percentiles, alongside stand and management fields such as planting year and known felling history (Table~\ref{tab:data-fields}). 
Figure~\ref{fig:data-study-area} shows the study-area topography, the spatial coverage of both cohorts, and an example of standard and boundary-clipped plots.

\begin{figure}[!htbp]
\includegraphics[
    width=1.05\textwidth,
    height=0.48\textheight,
    keepaspectratio
]{figures/fig_data/aberfoyle_area.png}
\caption{\small{Left:Aberfoyle study area. Middle: cohort coverage, Right: plot-level detail}}
\label{fig:data-study-area}
\end{figure}

Due to a LiDAR coverage increase of $37\text{~km}^2$ between 2006 and 2008, two nested cohorts were created (Figure~\ref{fig:cohort_coverage_funnel}): a \emph{4-survey} ($\in \{2008, 2012, 2021, 2023\}$) and a \emph{6-survey} ($\in \{2002, 2006, 2008, 2012, 2021, 2023\}$). The former has broad spatial coverage, while the latter extends temporal trajectories. Sequential filters required complete survey presence, pure Sitka spruce, planting between 1850 and the survey year, ages $0\text{--}200$ years, heights $0\text{--}60\text{~m}$ (ecological max; \cite{mason2011Sitka_C1_SS}), and age $\ge 30$ in 2023 to avoid unreliable young-stand heights (Table~\ref{tab:data-cleaning-funnel} \cite{suarezminguez2008Practical_C1_SS}).

\begin{table}[!htbp]
\centering
\scriptsize
\begin{tabular}{lrrrr}
\toprule
& \multicolumn{2}{c}{4-survey} & \multicolumn{2}{c}{6-survey} \\
\cmidrule(lr){2-3}\cmidrule(lr){4-5}
Stage & Plots & Removed & Plots & Removed \\
\midrule
0. Raw dataset
& 230,359 & -- & 230,359 & -- \\
1. Cohort years present
& 230,345 & 0.01\% & 230,359 & 0.00\% \\
2. Complete survey coverage
& 112,782 & 51.04\% & 18,781 & 91.85\% \\
3. Sitka spruce only
& 72,395 & 35.81\% & 13,897 & 26.01\% \\
4. Valid planting year
& 71,766 & 0.87\% & 13,897 & 0.00\% \\
5. Plausible age (0--200 years)
& 71,766 & 0.00\% & 13,897 & 0.00\% \\
6. Plausible height (0--60 m)
& 71,766 & 0.00\% & 13,897 & 0.00\% \\
7. Mature stands (age $\geq 30$ in 2023)
& \textbf{58,112} & 19.03\%
& \textbf{13,769} & 0.92\% \\
\bottomrule
\end{tabular}
\caption{\small{Plots left after each filter. Percentages show the share removed from the previous row, rather than from the raw dataset. Figure~\ref{fig:data-study-area} shows the final spatial coverage.}}
\label{tab:data-cleaning-funnel}
\end{table}

%------------------------------------------------
\section{Target variable and original stand fields}
\label{sec:data-response}


The target is the plot-level 95th percentile ALS return height above ground (\texttt{elev\_\allowbreak percentile\_\allowbreak 95th} per survey. Records also supply planting year, management history, and Forest Research yield and volume estimates (Table~\ref{tab:data-fields}). To prevent data leakage, any fields derived directly from the height target were excluded from model inputs.

\begin{table}[!htbp]
\scriptsize

\setlength{\tabcolsep}{3pt}
\renewcommand{\arraystretch}{0.95}
\begin{tabular}{
    >{\raggedright\arraybackslash}p{2.6cm}
    >{\raggedright\arraybackslash}p{7.0cm}
    >{\raggedright\arraybackslash}p{3.2cm}
}
\hline
Field & Meaning or derivation & Role in this study \\
\hline
\texttt{elev\_\allowbreak percentile\_\allowbreak 95th} & 95th percentile of ALS return heights above ground, used without correction & Response variable \\
\texttt{Top\_Height95} & $=\texttt{elev\_percentile\_95th}\times1.1$; Forest Research's own corrected height estimate & Retained for reporting but not modelled \\
\texttt{plyr}, \texttt{Age} & Planting year; \texttt{Age} = survey year $-$ \texttt{plyr} & Input feature; used for data cleaning \\
\texttt{yldc} & An existing Forest Research yield-class field, not calculated from this project's LiDAR-derived height response & Excluded from prediction; used as a comparison variable \\
\texttt{whcl} & Ordinal Forest Research windthrow hazard class, 0-6 & Management descriptor \\
\texttt{CanopyCover} & Canopy-cover fraction from the same ALS survey and processing step as the height response & Input feature \\
\texttt{Thin}, \texttt{last\_thinn}, \texttt{next\_thin\_} & Thinning history; \texttt{Thin} is binary (1 if the stand has been thinned, 0 if never thinned); \texttt{last\_thinn}/\texttt{next\_thin\_} record the last and next thinning years & Input feature \\
\texttt{Vol95}, \texttt{area}, \texttt{g1}, \texttt{g3} & Forest Research volume estimate, calculated from height, plot area and canopy cover & Used for disturbance screening \\
\texttt{GYCspec95}, \texttt{p1-p5} & General Yield Class, calculated directly from a plot's own measured height (\texttt{Top\_Height95}) & Excluded from prediction because its calculation
incorporates the height response, creating target leakage \\
\hline
\end{tabular}
\caption{\small{Target, stand and management fields provided by Forest Research.}}
\label{tab:data-fields}
\end{table}
%------------------------------------------------
\section{Derived terrain, wind, climate and soil variables}
\label{sec:data-environmental}

Each plot was matched with external terrain, wind, climate, and soil variables (Table~\ref{tab:data-env-sources}). Appendix~\ref{tab:data-env-variables_detail_app} defines each variable.

\begin{table}[!htbp]
\centering
\scriptsize
\setlength{\tabcolsep}{4pt}
\renewcommand{\arraystretch}{1.05}

\begin{tabular}{
    >{\raggedright\arraybackslash}p{2.0cm}
    >{\raggedright\arraybackslash}p{4.2cm}
    >{\raggedright\arraybackslash}p{2.3cm}
    >{\raggedright\arraybackslash}p{4.0cm}
}
\hline
Data family & Source and resolution & Time basis & Variables used \\
\hline

Terrain
& OS Terrain 50, 50\,m \cite{ordnancesurveyTerrain50_C3_LID}
& Static
& Elevation, slope, aspect, northness and eastness \\

Soil
& CEH Natural Capital, 50\,m \cite{feeney2026CEHNaturalCapital_C3_LID}; SoilGrids v2.0, 250\,m \cite{poggio2021SoilGrids_C3_LID}
& Static
& Pedotope, drainage, texture and soil pH \\

Wind exposure
& Global Wind Atlas, 250\,m \cite{davis2023GlobalWindAtlas_C3_LID}
& 2006--2015 climate
& Long-term wind speed and exposure \\

Observed wind
& MIDAS Open weather stations \cite{metoffice2019MIDASOpen_C3_LID}
& Survey intervals
& Mean and extreme wind conditions \\

Long-term climate
& CHELSA v2.1, 1\,km \cite{karger2017CHELSA_C3_LID}
& 1981--2010 climate
& Growing degree days and precipitation \\

Matched climate
& HadUK-Grid, 1\,km \cite{hollis2019HadUKGrid_C3_LID}
& Cohort years
& Temperature and frost \\

Spatial context
& OS Open Roads and Open Rivers, vector data
\cite{ordnancesurveyOpenRoads_C3_LID,ordnancesurveyOpenRivers_C3_LID}
& Static
& Distance to roads, rivers and forest edges \\

\hline
\end{tabular}

\caption{\small{External data sources used in the analysis. Values were matched to plot centres. Full source, resolution, coverage and extraction details are given in Appendix~\ref{app:data-sources}.}}
\label{tab:data-env-sources}
\end{table}

The Global Wind Atlas gives long-term modelled wind, TOPEX measures terrain shelter, and MIDAS records wind at three weather stations. CHELSA describes long-term climate from 1981 to 2010 \cite{karger2017CHELSA_C3_LID}. HadUK-Grid temperature and frost values were averaged across each cohort's survey years \cite{hollis2019HadUKGrid_C3_LID}.


%=======================================================================================
%=======================================================================================
\chapter{Methodology and Experimental Design} \label{ch:methodology}
This chapter outlines how the raw data was utilised were used to answer the three research questions. It defines the splits, feature sets, models, targets and evaluation procedures. The same held-out observations were used when models were compared.

\paragraph*{Experimental Framework:}RQ1 tests predictive accuracy. RQ2a tests when environmental conditioning corrects departure
from a shared growth curve. RQ2b attributes persistent departure from that curve. RQ3 attributes
plot-specific asymptote deviation and tests whether the associations vary over space.
\begin{table}[!ht]
\centering
\footnotesize
\renewcommand{\arraystretch}{1.15}

\begin{tabular}{
@{}
>{\raggedright\arraybackslash}p{0.9cm}
>{\raggedright\arraybackslash}p{3.5cm}
>{\raggedright\arraybackslash}p{3.2cm}
>{\raggedright\arraybackslash}p{2.8cm}
@{}
}
\toprule
RQ & Unit and target & Models or comparison & Main evaluation \\
\midrule
RQ1 & Plot--survey row; 95th-percentile top height & CR, linear regression, random forest,
XGBoost, DNN, PINN and PINN-$k$ & Five-fold spatial block CV; temporal and plot-level checks \\
RQ2a & Plot--survey row; error reduction from the fold-matched CR curve & Conditioned
DNN, PINN, PINN-$k$ and XGBoost; no-env and permutation controls & Five-fold spatial block
CV; CR-error quartiles \\
RQ2b & Plot; mean pooled-CR residual & Elastic Net, XGBoost and two-stage
mixed-effects model & Five-fold spatial block CV; coefficients, SHAP and residual Moran's $I$ \\
RQ3 & Plot; difference between fitted and yield-class-implied asymptotes & Elastic Net, XGBoost
and GNNWR & Five-fold spatial block CV; attribution, Moran's $I$ and target audit \\
\bottomrule
\end{tabular}
\caption{Units, targets and evaluation used for each research question.}
\label{tab:rq-framework}
\end{table}

\begin{figure}[!htbp]
\centering
\includegraphics[width=\textwidth]{figures/fig_methodology/method_fig_pipeline-overview.png}
\caption{\small{Analysis pipeline. The two cohorts feed shared spatial partitions and separate RQ-specific feature sets. RQ1 predicts top height, RQ2 examines departure from a shared curve, and RQ3 examines plot-specific asymptote deviation.}}
\label{fig:pipeline-overview}
\end{figure}

\section{Data partitions and leakage control}
The four-survey cohort was used as the main dataset, while the smaller six-survey cohort provided a nested sensitivity analysis with longer observation histories. The unit of analysis depended on the research question. RQ1 and RQ2a used plot--survey observations, whereas RQ2b and RQ3 aggregated the data to one observation per plot.

\subsection*{Evaluation splits}
Spatial-block cross-validation was the primary evaluation design. Whole forest compartments were assigned to folds, preventing the same compartment from appearing in both training and testing \cite{roberts2017Cross_C5_EVAL}. Training plots within 60\,m of a validation or test plot were also removed, while all held-out plots were retained. Five folds were balanced approximately by row count. In each rotation, one fold was used for testing, the following fold for validation and the remaining three for training. Each compartment was therefore tested once and predictions from the five test folds were pooled before calculating the overall metrics \cite{ploton2020Spatial_C5_EVAL}. 

A random plot-level split assigned plots to training, validation and testing in a 60/20/20 ratio. This provided an easier reference without spatial separation. The temporal split instead separated observations by survey year \cite{lynch2025Digital_C4_NN}. For the 4-survey cohort, 2008 and 2012 were used for training, 2021 for validation and 2023 for testing. For the 6-survey cohort, 2002--2012 were used for training, followed by validation on 2021 and testing on 2023. Random assignments used seed~42.

\begin{figure}[!htbp]
\includegraphics[width=\textwidth]{figures/fig_methodology/method_fig-splits-combined_draft.png}
\caption{Evaluation data-splitting strategies: (A) random plot-level split, (B) compartment-level spatial-block split with a 60\,m buffer, and (C) temporal holdout timeline by survey year.}
\label{fig:splits-combined}
\end{figure}


Continuous predictors were standardised using parameters estimated from the training partition. The same transformation was applied to validation and test data. Age was scaled separately so that the PINN derivative could be converted to metres per year. Categorical soil classes were one-hot encoded.Compartment, sub-compartment and block identifiers were excluded from the model inputs. Two neighbour-height variables were also excluded because they had been calculated from nearby plots' observed heights before partitioning and could therefore leak held-out outcomes.

%------------------------------------------------
\section{Environmental feature sets}

Because RQ1, RQ2b, and RQ3 target different response variables, feature sets are constructed separately for each question using a shared four-step protocol:

\begin{enumerate}
\setlength\itemsep{0pt}
\item \textbf{Coverage screening:} Variables with insufficient spatial coverage are removed prior to ranking.
\item \textbf{Collinearity filtering:} For collinear pairs with Spearman $|\rho| \ge 0.95$, one representative is retained, prioritising direct external measurements over derived indices.
\item \textbf{Composite ranking:} Remaining candidates are sorted by their average rank across univariate Spearman $|\rho|$, permutation feature importance, and drop-column loss in held-out $R^2$.
\item \textbf{Multicollinearity pruning:} Features are pruned iteratively until all variance inflation factors (VIF) fall below 5 \cite{obrien2007Caution_C5_EVAL}. Protected stand variables remain fixed in the matrix throughout this process to constrain the VIF of competing predictors.
\end{enumerate}

Four nested feature sets are defined from these steps. Set~1 contains only the protected baseline variables: canopy cover, time since thinning, its missingness indicator, and recent thinning. Set~2 adds the ten highest-ranked candidates. Set~3 expands on Set~2 by adding the top half of terrain and wind features, while Set~4 adds the top half of climate, soil, and spatial context features to Set~3. Appendix~\ref{app:feature-sets} lists complete variable assignments and VIF diagnostics.

Feature selection is performed once on the designated training partition rather than re-run across individual folds, which ensures all models train on identical feature subsets, meaning fold-to-fold variance reflects true model stability and spatial partitioning rather than shifting predictor selection.

%------------------------------------------------
\section{RQ1: Top-height prediction}
\label{sec:rq1-method}

RQ1 compares a biological curve, standard tabular models and neural models. The headline
comparison uses Set~3 and the same five spatial folds. The other feature sets test whether adding
environmental inputs changes the result.

\subsection{Chapman--Richards (CR) baseline}
\begin{equation}
H(a) = y_{\max}\left(1 - e^{-ka}\right)^p,
\label{eq:cr}
\end{equation}
where $H(a)$ is top height at age $a$, $y_{\max}$ is the asymptotic height ceiling, $k$ is the growth rate parameter, and $p$ governs curve shape. The model is fitted using non-linear least squares across five initial starting points, retaining the solution with the lowest sum of squared errors. Parameters are calibrated exclusively on the training partition of each split and fold, then frozen. This prevents validation and test observations from influencing the baseline fit or the PINN loss terms derived from it.

\subsection*{Conventional models}
Three classical baselines represent different levels of complexity: ordinary linear regression, a bagged random forest, and gradient-boosted decision trees via XGBoost \cite{chen2016XGBoost_C4_NN}. All conventional models use the same RQ1 feature sets and cross-validation folds.

For RQ1, XGBoost hyperparameters were tuned specifically on the top-height target. Evaluating 27 candidate configurations across five folds and two cohorts produced 270 fits. Model selection relied strictly on validation $R^2$, keeping test sets untouched until final evaluation. Appendix~\ref{app:hyperparameters} lists the full search grid and chosen values. For RQ2b and RQ3, attribution models use a fixed configuration instead: tree depth 4, learning rate 0.04, and up to 500 boosting rounds with early stopping, without further retuning.

\subsection{DNN and physics-informed models}
All neural architectures share a three-layer prediction trunk with 128 units per layer and Leaky-ReLU activations. Training uses the Adam optimiser with a base learning rate of $10^{-4}$, weight decay of $10^{-5}$, gradient clipping at 1.0, batch size 256, and early stopping with a patience of 40 epochs (up to 500 epochs total). A ReduceLROnPlateau scheduler (decay factor 0.8, patience 15 epochs) lowers the learning rate when validation loss stalls, remaining active across most of the optimisation trajectory. The standard DNN minimises mean squared error:
\begin{equation}
\mathcal{L}_{\mathrm{data}}=\frac{1}{N}\sum_{i=1}^{N}
\left(\hat H_i-H_i\right)^2.
\label{eq:data-loss}
\end{equation}

The environment-conditioned DNN concatenates the environmental variables to the main input. The PINN variants use the same prediction trunk, but environmental variables enter small auxiliary networks. The architectures therefore share the same core capacity, while the environmental information is routed differently.

Routing environment data through side networks tests a specific growth-ceiling effect rather than a broad, mixed one. Feeding terrain directly into the main trunk would blend site effects into general sensitivity, making them hard to separate from the standard DNN baseline (Section~\ref{sec:results-rq1-cr-reduction}).

The analytical derivative of the CR curve is:
\begin{equation}
H'(a)=y_{\max}pk e^{-ka}\left(1-e^{-ka}\right)^{p-1}.
\label{eq:cr-derivative}
\end{equation}
The physics loss penalises discrepancies between this theoretical growth rate and the network derivative computed via automatic differentiation:
\begin{equation}
\mathcal{L}_{\mathrm{phys}}=\frac{1}{N}\sum_{i=1}^{N}
\left(\frac{\partial\hat H_i}{\partial a_i}-H'(a_i)\right)^2.
\label{eq:physics-loss}
\end{equation}
To maintain consistent units ($\mathrm{m\,yr^{-1}}$), the network derivative is divided by the training-set age standard deviation before computing the loss.

The trajectory loss uses consecutive training surveys from the same plot:
\begin{equation}
\mathcal{L}_{\mathrm{traj}}=\frac{1}{M}\sum_{j=1}^{M}
\left(\frac{\hat H_{j,2}-\hat H_{j,1}}{\Delta a_j}-H'(\bar a_j)\right)^2,
\label{eq:trajectory-loss}
\end{equation}
where $\bar a_j$ is the midpoint age. Pairs are formed only when both rows belong to the training
partition. The total loss is
\begin{equation}
\mathcal{L}=\mathcal{L}_{\mathrm{data}}+
\lambda_{\mathrm{phys}}\mathcal{L}_{\mathrm{phys}}+
\lambda_{\mathrm{traj}}\mathcal{L}_{\mathrm{traj}}+
10^{-5}\sum_{\theta}|\theta|.
\label{eq:total-loss}
\end{equation}
The default uses $\lambda_{\mathrm{phys}}=\lambda_{\mathrm{traj}}=1$. Joint weights 0, 1 and 2
are compared in the ablation. Weight 0 is the physics-disabled control.

The first conditioned PINN changes the asymptote:
\begin{equation}
y_{\max}^{(i)}=y_{\max}+f_y(\mathbf{x}_i).
\label{eq:ymax-adjust}
\end{equation}
PINN-$k$ also changes the rate while keeping it positive:
\begin{equation}
k^{(i)}=k\exp\!\left(f_k(\mathbf{x}_i)\right).
\label{eq:k-adjust}
\end{equation}
Both auxiliary networks have one 16-unit hidden layer and start near zero. The shape parameter
$p$ remains fixed in every variant.

The trunk's prediction never evaluates Equation~\ref{eq:cr} with $y_{\max}^{(i)}$ or $k^{(i)}$;
these enter only the physics and trajectory losses (Equations~\ref{eq:physics-loss}--
\ref{eq:trajectory-loss}). Environmental information therefore reaches the trunk only indirectly,
through those loss gradients, never as a direct input to the prediction itself.

%------------------------------------------------
\section{RQ2a: Environmental correction of curve departure}

RQ2a evaluates whether environmental predictors explain departures from standard yield curves using held-out predictions from RQ1. For each test observation $i$, the absolute error reduction relative to the fold-matched Chapman--Richards baseline is
\begin{equation}
d_i=\left|H_i-\hat H_i^{\mathrm{CR}}\right|-
\left|H_i-\hat H_i^{\mathrm{model}}\right|.
\label{eq:residual-reduction}
\end{equation}
A positive value indicates that the model improves on the baseline prediction. The analysis reports mean $d_i$, its percentage of mean Chapman--Richards error, and the percentage of observations where $d_i>0$.

Within each held-out fold, test observations are ranked by absolute baseline error and divided into quartiles (Q1 to Q4). Q1 contains observations where the Chapman--Richards curve fits best, while Q4 contains those with the largest departures. Mean error reduction and the percentage of improved plots are calculated for each quartile. This is repeated for the environment-conditioned DNN, PINN and PINN-$k$ across Sets~2--4 and both cohorts, and for XGBoost on Set~3 and both cohorts. The Set~3 DNN is evaluated across five training seeds to measure parameter initialisation spread. Other RQ2a configurations use seed 42, meaning their reported spread reflects fold-to-fold variation rather than seed variance.

Two controls, both run on XGBoost with Set~3 only, isolate the source of any improvement. First, XGBoost is evaluated with and without environmental inputs. Second, the environmental columns are permuted using a single shared row permutation across every row of a held-out fold (train, validation, and test alike), and XGBoost is refitted from scratch on the permuted data. This tests whether the performance gain requires real environment-plot linkage, avoiding the confound of evaluating an already-trained model on out-of-distribution inputs. The same error-reduction calculation is then repeated on the refitted model's held-out predictions. For the DNN and XGBoost (both Set~3, both cohorts), the spatial pattern of $d_i$ is assessed with Moran's $I$ using the same weights as the residual analysis.

%------------------------------------------------
\section{RQ2b: Persistent departure from the shared curve}

RQ2b focuses on the four-survey cohort, as the six-survey cohort contains too few compartments for compartment-level spatial attribution. For each plot-survey observation, the deviation from expected height is:
\begin{equation}
r_{it}=H_{it}-H(a_{it}),
\end{equation}
yielding the time-averaged plot target:
\begin{equation}
\bar r_i=\frac{1}{T_i}\sum_{t=1}^{T_i}r_{it}.
\label{eq:mean-cr-residual}
\end{equation}
Positive values identify stands that consistently outperform the pooled Chapman--Richards curve. Stand age is excluded from the candidate predictor matrix because it already defines the baseline expectation.

Elastic Net and XGBoost are evaluated under five-fold compartment cross-validation. Elastic Net tunes its penalty parameter and $\ell_1$ mixing ratio through nested cross-validation on each training partition. XGBoost uses the fixed hyperparameters specified in Section~\ref{sec:rq1-method}. Model performance is measured using pooled out-of-fold $R^2$.

To account for hierarchical spatial grouping, a two-stage mixed-effects model is also estimated. Stage~1 computes $\bar r_i$ from the non-linear baseline curve. Stage~2 fits:
\begin{equation}
\bar r_i=\boldsymbol{\beta}^{\top}\mathbf{x}_i+u_{c[i]}+\varepsilon_i,
\label{eq:mixed-effects}
\end{equation}
where $u_{c[i]}\sim\mathcal{N}(0,\sigma_u^2)$ represents a random intercept for each compartment, without random slopes. Fixed-effect coefficients are estimated using restricted maximum likelihood within each training fold. Spatial variance explained is calculated as the proportional reduction in compartment random-intercept variance relative to an intercept-only baseline. Parameters are reported as the mean $\pm$ sample standard deviation across folds.

Four diagnostic checks aid interpretation. First, Set~4 is re-estimated without canopy cover. Second, Pearson correlations are calculated between canopy cover and each target variable. Third, SHAP values are recomputed without canopy cover to determine whether correlated predictors absorb its importance. Fourth, between-compartment variance shares are computed for slope and TOPEX to resolve differences between the mixed model and Elastic Net.

Together, Elastic Net coefficients, mixed-effects fixed effects, and XGBoost mean absolute SHAP values provide three distinct views of growth attribution. SHAP calculations are performed post-estimation and do not influence feature selection. Finally, Moran's $I$ is computed on the pooled held-out residuals from Elastic Net and XGBoost to check for remaining spatial autocorrelation.

%------------------------------------------------
\section{RQ3: Plot-specific asymptote deviation}
\label{sec:rq3-method}

RQ3 estimates a distinct asymptotic height ceiling for each plot while fixing the Forestry Commission rate and shape parameters. For plot $i$ at survey time $t$, let
\begin{equation}
g_{it}=\left(1-e^{-p_{4,i}a_{it}}\right)^{p_{5,i}}.
\end{equation}
Fitting the single-parameter growth curve through the origin by ordinary least squares yields
\begin{equation}
\hat y_{\max,i}=\frac{\sum_t H_{it}g_{it}}{\sum_t g_{it}^2}.
\label{eq:plot-ymax-fit}
\end{equation}
The expected asymptote implied by standard yield class is calculated directly from the regional Forestry Commission parameters:
\begin{equation}
y_{\max,i}^{\mathrm{yldc}}=\mathrm{yldc}_i p_{2,i}+p_{1,i}+2p_{3,i}.
\end{equation}
The final attribution target is the deviation between these two ceilings:
\begin{equation}
\Delta y_{\max,i}=\hat y_{\max,i}-y_{\max,i}^{\mathrm{yldc}}.
\label{eq:local-ymax-diff}
\end{equation}
This formulation keeps the target strictly free of environmental covariates. Fixing $p_4$ and $p_5$ avoids over-parameterising curves from only four to six time points, though it absorbs any local shape mismatch into the fitted asymptote.

Elastic Net and XGBoost models are trained on $\Delta y_{\max}$ using feature Sets~2--4. To model spatially varying associations, Geographically Neural Network-Weighted Regression (GNNWR) is also estimated \cite{du2020Geographically_C6_GWR}:
\begin{equation}
\hat y_i=\beta_0(s_i)+\sum_j\beta_j(s_i)x_{ij}.
\label{eq:gnnwr}
\end{equation}
For each target location $s_i$, the spatial weighting sub-network receives a MinMax-scaled vector of Euclidean distances to all training reference plots. The network predicts multiplicative scaling factors that adjust a frozen, global ordinary least squares baseline for the intercept and each slope parameter. All runs use the full, uncapped training reference set without predefined distance kernels or fixed bandwidths. Predictions for held-out validation plots reference these same training points, evaluated across the shared five spatial folds.

Predictor importance is assessed using both tree split gain and post-hoc SHAP values. XGBoost gain sums the training loss reduction across all tree splits involving a feature, reflecting its utility during model construction. In contrast, mean absolute SHAP values quantify each feature's average additive contribution to final predictions \cite{lundberg2017unified}. RQ3 SHAP values are computed out of fold across each test partition and pooled before interpretation.


\subsection*{Target and outlier audit}
To prevent flawed time series from corrupting the analysis, a plot is excluded from the curve-fitting population if any consecutive survey interval meets either of two filtering rules:
\begin{itemize}
\setlength\itemsep{2pt}
\item \textbf{Clearfell pattern:} Top height falls by at least 25\%, canopy cover drops by at least 70\%, and standing volume drops by at least 70\%.
\item \textbf{Measurement inconsistency:} Top height falls by at least 25\%, while canopy cover drops by no more than 10\% and standing volume does not decrease.
\end{itemize}
Intervals exhibiting large height drops that match neither criterion are classified as ambiguous and kept in the dataset.

Unusually low or high fitted asymptotes are screened separately using Tukey's $1.5\,\mathrm{IQR}$ fences \cite{tukey1977Exploratory}:
\begin{equation}
\left[Q_1 - 1.5\,\mathrm{IQR},\; Q_3 + 1.5\,\mathrm{IQR}\right],
\label{eq:tukey-fences}
\end{equation}
calculated independently within each cohort. Similar non-parametric thresholds have been used to identify parameter estimation errors in non-linear growth modelling \cite{markovic2014NewApproach}. Because this fence captures distributional extremes rather than demonstrably erroneous records, flagged plots remain in the primary analysis. A sensitivity check recalculates pooled held-out metrics after removing these observations from saved predictions, without refitting the models.

The audit evaluates whether specific plots recur within the worst 1\% of residuals across models. It compares their within-plot residual range, edge distance to compartment boundaries, spatial clustering, and topographic wind shelter against the wider sample. To test a known survey-boundary issue, the worst-fitting survey year is identified, and each flagged plot is refitted after excluding its 2008 observation. These diagnostics assess whether errors originate from target construction rather than genuine ecological factors.

\section{Evaluation and robustness}

For residual $e_i=H_i-\hat H_i$, the main metrics are
\begin{align}
\mathrm{RMSE}&=\sqrt{\frac{1}{N}\sum_i e_i^2}, &
\mathrm{MAE}&=\frac{1}{N}\sum_i|e_i|, \\
R^2&=1-\frac{\sum_i e_i^2}{\sum_i(H_i-\bar H)^2}, &
\mathrm{Bias}&=\frac{1}{N}\sum_i e_i.
\end{align}
Positive bias denotes under-prediction. Each table states whether it reports a fold mean or a
pooled metric. Fold summaries are the mean $\pm$ sample SD across five folds. Pooled metrics are
calculated once from all out-of-fold predictions. These are not interchangeable.

For the four-survey spatial analyses, pooled $R^2$ 95\% confidence intervals use 1,000 cluster bootstrap resamples of whole compartments. The 2.5th and 97.5th percentiles form the interval. This preserves spatial clustering better than resampling individual rows. The smaller six-survey RQ3 results are reported as point estimates; GNNWR is also repeated over three split seeds to test sign stability.

A semivariogram measures how quickly residuals become less similar as distance between plots increases; its range is the distance beyond which they are no longer spatially correlated. Each model's range is fitted separately from its own residuals, searched up to 5000m by default (widened to 10,000m where a fit did not resolve within that window); two models' Moran's $I$ are not necessarily tested at the same distance.

Global Moran's $I$ is calculated from held-out residuals using distance-band weights, with the threshold set by this semivariogram range, and significance assessed via 999 permutations \ref{app} \app{add semivariogram ranges used to app} 

\begin{equation}
I=\frac{N}{W}
\frac{\sum_i\sum_j w_{ij}(e_i-\bar e)(e_j-\bar e)}
{\sum_i(e_i-\bar e)^2}.
\label{eq:morans-i}
\end{equation}
It is used after fitting, not for feature selection. 
Three outcomes occur: the range resolves within the search window (Moran's $I$ reported directly); the semivariance stays flat throughout the window (no spatial structure detected, Moran's $I$ not computed); or the semivariance is still rising at the window's edge (Moran's $I$ computed at that capped distance, reported as provisional, not equivalent to a resolved result). The same rule is applied to RQ2a error reduction, RQ2b residuals and the Elastic Net, XGBoost and GNNWR residuals in RQ3.

Neural training variation is assessed over five seeds for the selected RQ1 models. XGBoost and the
DNN are also compared fold by fold on identical test folds. The physics-weight ablation tests
weights 0, 1 and 2. For PINN-$k$, parameter separability is assessed by the pooled correlation
between learned $y_{\max}$ and $k$. Its confidence interval uses the same 2,000 compartment
bootstrap resamples. These checks distinguish predictive accuracy from parameter stability.

\section{Reproducibility}

Split assignments, feature lists and model settings are saved with each run. Seed 42 is used unless
a stated reseed is part of the analysis. Full feature membership, hyperparameter searches,
secondary split results and detailed outlier tables are retained in the Appendix. Methodological
limitations and their effect on interpretation are evaluated after the results, before the final
conclusions.

%=======================================================================================
%=======================================================================================
\chapter{Results and Discussion} \label{ch:results}

This chapter covers three research questions. RQ1 tests top-height prediction using spatial-block cross-validation, followed by temporal and diagnostic checks. Set 3 is the baseline for RQ1, and it gives all models shared terrain and wind data. RQ2 tracks persistent plot departures from the Chapman–Richards curve, and RQ3 tests if global or spatial models explain plot-level asymptote shifts. Other feature sets are in \app{}

\section{RQ1: prediction of top heights} 

A tuned XGBoost model gives the most accurate top-height predictions of every model
tested, on both cohorts, under the spatial-block evaluation. CR guidance does not recover this
gap at any tested loss weight (\textit{lambda}), but it is useful for parameter identifiability (estimating both $y\_max$ and growth rate $k$. 

\subsection{Results: top-height prediction:}
\begin{table}[!htbp]
\centering
\scriptsize
\setlength{\tabcolsep}{4pt}
\renewcommand{\arraystretch}{1.2}
\begin{tabular}{lrrrrrr}
\toprule
& \multicolumn{3}{c}{Four-survey} & \multicolumn{3}{c}{Six-survey} \\
\cmidrule(lr){2-4} \cmidrule(lr){5-7}
Model & RMSE${\downarrow}$ & MAE${\downarrow}$ & $R^2{\uparrow}$ & RMSE${\downarrow}$ & MAE${\downarrow}$ & $R^2{\uparrow}$ \\
\midrule
Linear & $5.168{\pm}0.361$ & $4.056{\pm}0.343$ & $0.580{\pm}0.026$ & $4.877{\pm}0.643$ & $3.619{\pm}0.426$ & $0.509{\pm}0.067$ \\
RF & $4.794{\pm}0.187$ & $3.611{\pm}0.142$ & $0.638{\pm}0.011$ & $4.062{\pm}0.249$ & $3.010{\pm}0.185$ & $0.656{\pm}0.046$ \\
XGBoost (tuned) & $\mathbf{4.548{\pm}0.137}$ & $\mathbf{3.431{\pm}0.123}$ & $\mathbf{0.674{\pm}0.009}$ & $\mathbf{3.656{\pm}0.232}$ & $\mathbf{2.669{\pm}0.183}$ & $\mathbf{0.722{\pm}0.031}$ \\
DNN & $4.681{\pm}0.230$ & $3.533{\pm}0.201$ & $0.655{\pm}0.016$ & $3.903{\pm}0.316$ & $2.867{\pm}0.275$ & $0.684{\pm}0.031$ \\
PINN & $5.206{\pm}0.399$ & $4.063{\pm}0.407$ & $0.573{\pm}0.036$ & $3.886{\pm}0.290$ & $2.819{\pm}0.196$ & $0.686{\pm}0.038$ \\
PINN-$k$ & $5.197{\pm}0.402$ & $4.046{\pm}0.421$ & $0.575{\pm}0.037$ & $3.895{\pm}0.282$ & $2.816{\pm}0.196$ & $0.684{\pm}0.038$ \\
\bottomrule
\end{tabular}
\caption{Model comparison for top-height prediction, Set3, spatial\_block\_kfold,
per-fold mean $\pm$ SD across five folds. 4survey n=232{,}292 row-years (58{,}073
plots); 6survey n=82{,}614 row-years (13{,}769 plots). Table shows pooled mean $\pm$ SD across all folds, and margin reflects the full split rather than single folds.}
\label{tab:results-rq1}
\end{table}

\begin{table}[!htbp]
\centering
\small
\renewcommand{\arraystretch}{1.2}
\setlength{\tabcolsep}{4pt}
\begin{tabular}{llrrr}
\toprule
Cohort & Model & Spatial-block $R^2$ & Temporal $R^2$ & Absolute change \\
\midrule
4survey & DNN & 0.634 & 0.544 & $-0.090$ \\
4survey & PINN & 0.584 & 0.290 & $-0.294$ \\
4survey & PINN$_k$ & 0.588 & 0.289 & $-0.299$ \\
6survey & DNN & 0.722 & 0.317 & $-0.405$ \\
6survey & PINN & 0.731 & 0.161 & $-0.570$ \\
6survey & PINN$_k$ & 0.730 & 0.190 & $-0.540$ \\
\bottomrule
\end{tabular}
\caption{Spatial-block vs.\ temporal-holdout $R^2$, by model and cohort. Training uses survey data
up to 2012, spatial-block validation uses 2021 data, temporal holdout tests on 2023 data. Only
$R^2$ was computed for this check; RMSE and MAE are not available from this source
\fc{(re-run if RMSE/MAE are needed for the final draft)}.}
\label{tab:results-rq1-temporal}
\end{table}

\subsection{Discussion}
XGBoost outperformed the DNN in both cohorts. The advantage was larger over the physics-guided
models in the four-survey cohort. A fold-by-fold check confirms it: XGBoost won 9 of the 10 individual folds tested across both cohorts (all 5 on 6survey; 4 of 5 on 4survey, losing one fold by a small margin, $-0.007$ $R^2$). Tests to understand why XGBoost was outperforming the neural models were: 

\begin{itemize}
\item \textbf{Feature interactions (insufficient explanation):} XGBoost allocates 27\% of its SHAP prediction weight to interactions (primarily Age crossed with terrain and stand structure). However, DNN already captures 80\% of the nonlinear gain above the linear baseline ($R^2$ 0.580 $\to$ 0.655 out of a possible 0.674; Table~\ref{tab:results-rq1}), and Random Forest underperforms the DNN ($R^2 = 0.638$). Interaction capacity alone therefore does not explain the margin.
\item \textbf{Compartment sample size (ruled out):} Subsampling the 4-survey cohort down to 48 compartments (matching the 6-survey count) did not widen XGBoost's lead; the $R^2$ difference fluctuated between $-0.085$ and $+0.029$ across random subsets, compared to $+0.024$ on full data.
\item \textbf{Architecture-size under-tuning (ruled out):} Sweeps across five network sizes showed no consistent gains for DNN, while PINN and PINN$_k$ showed even lower sensitivity to model capacity.
\item \textbf{Terrain extrapolation in held-out compartments (ruled out):} On the 4.1\% of 4-survey test plots with out-of-range terrain values, XGBoost degraded more relative to its in-range error than DNN did (error ratio 1.312 vs.\ 1.243; correlation with out-of-range distance $\rho=0.074$ vs.\ $0.059$).
\item \textbf{Training hyperparameter under-tuning (ruled out):} A 48-configuration sweep over dropout and learning rate yielded no meaningful gains for neural models under 5-fold cross-validation (within $\pm0.008$ $R^2$ of default baselines). Similarly, XGBoost's 27-configuration tuning shifted $R^2$ by only 0.002--0.004, confirming that neither model family had substantial untapped hyperparameter headroom.
\end{itemize}

The most likely remaining account is that XGBoost, as a boosted-tree model, simply suits this tabular prediction task better than this DNN architecture does, consistent with known theory that boosted trees remain strong on heterogeneous tabular data \cite{shwartzziv2021Tabular_C4_NN}.  

\paragraph*{Cohort and Temporal sensitivity:}
The 6-survey cohort represents a distinct slice of the forest rather than simply a smaller sample. It covers fewer compartments (48 vs.\ 296), occupies a lower and more uniform elevation ($102.3\pm47.9$~m, range 18--351~m, vs.\ $177.9\pm104.6$~m, range 18--561~m), and yields a growth curve with a lower ceiling and faster initial growth rate ($y_{\max}=44.46$~m, $k=0.0233$ vs.\ $y_{\max}=51.96$~m, $k=0.0103$). This environmental divergence explains why the 6-survey cohort behaves differently throughout this chapter.

All models suffered accuracy losses under the 2023 temporal test relative to spatial-block evaluation (Table~\ref{tab:results-rq1-temporal}), consistent with the challenge of projecting forward in time. However, physics guidance worsened this drop rather than preventing it. PINN and PINN$_k$ lost far more accuracy than the unconstrained DNN across both cohorts (4-survey: DNN $-0.090$ vs.\ PINN $-0.294$; 6-survey: DNN $-0.405$ vs.\ PINN $-0.570$). If the Chapman–Richards constraint aided out-of-period extrapolation, the PINN drop would be smaller; instead, it was substantially larger, marking a second clear limitation of physics guidance alongside its lack of spatial advantage. The larger decline in the 6-survey cohort is due to the homogeneity in this subset, where stopping training in 2012 and testing into 2023, means its harder to generalise. 

\paragraph*{Effect of the evaluation split}
Switching from a spatial-block to a random plot-level split \ref{fig:splits-combined}, sharply inflates the DNN's score (4-survey: $0.831$ vs.\ $0.634$; 6-survey: $0.844$ vs.\ $0.722$). Because adjacent plots share terrain and stand conditions, random splits leak spatial information from training into testing and overstate true performance. This data leak is why spatial-block splits are used as the main test across this chapter.

The PINN does not get this boost, and its scores stay almost the same across splits (0.590 vs.\ 0.584 on 4survey, 0.728 vs.\ 0.731 on 6survey). This happens because terrain data never enters the PINN's main trunk, so there is no local info to leak between train and test. A plain DNN with no terrain inputs (\textit{dnn\_noenv}) confirms this: it has the same small split gap as the PINN (0.657 vs.\ 0.633 on 4survey), and its spatial score (0.633) matches the regular DNN (0.634) almost exactly. This rules out physics rules and rules out a weak fit, proving that feeding terrain into the main trunk is what causes the leak.

\paragraph*{Physics-weight ablation} Raising the physics loss weight ($\lambda \in \{0, 1, 2\}$) made predictions worse every time, across both models and cohorts. Setting $\lambda=0$ raised $R^2$ by $0.05$ on 4survey (a real gain) and by $+0.010$ on 6survey (within noise). The physics loss directly hurts accuracy by fighting the data fit on shared weights. Dropping the physics loss cuts the 4survey DNN--PINN gap from $0.080$ down to $0.026$ (DNN $0.655$ vs.\ unconstrained PINN$_k$ $0.629$), and it closes the gap on 6survey ($0.719$ vs.\ $0.684$). 

The cause of the remaining $0.026$ gap on 4survey is unknown. One theoretical possibility: PINN$k$ carries extra parameters (the $y{\max}$/$k$ sub-networks) that exist in the model regardless of $\lambda$. If weight decay and the L1 penalty apply uniformly across all parameters, PINN$_k$'s effective regularisation is spread over a larger parameter count than the DNN's, even with the physics loss contributing zero gradient, this alone could produce a small accuracy gap unrelated to physics guidance itself. Equally plausible: this is ordinary run-to-run variation between two structurally different networks, with no special mechanism required. 

\paragraph*{Answer to RQ1.}
Tuned XGBoost gave the strongest spatially held-out predictions in both cohorts. The DNN remained competitive, but CR guidance did not improve predictive accuracy. This is architectural: terrain never reaches PINN or PINN$_k$'s predictions, at any weight, only the training loss. Random plot splitting gave an optimistic DNN estimate. The current $\lambda=0$ implementation does not support a controlled claim about parameter identifiability.






%=======================================================================================
%=======================================================================================
\chapter{Conclusion and Future Work} \label{ch:conc}

%=======================================================================================
%=======================================================================================
\bibliographystyle{plain}
\bibliography{mybibfile}
%=======================================================================================

\appendix
\chapter{Appendix}  \label{chap:appendix} 
% \TODO: create appendix secions and labels for each section. 



%=======================================================================================
\end{document}
%=======================================================================================

